\documentclass[11pt,a4paper]{article}
% Shared preamble of the FEBAUTO Racing manuals.
\usepackage[margin=2.2cm,headheight=16pt]{geometry}
\usepackage[T1]{fontenc}
\usepackage{lmodern}
\usepackage{xcolor}
\usepackage{graphicx}
\usepackage{booktabs}
\usepackage{tabularx}
\usepackage{listings}
\usepackage{fancyhdr}
\usepackage{enumitem}
\usepackage{tikz}
\usetikzlibrary{arrows.meta}
\usepackage[hidelinks,colorlinks=true,linkcolor=febblue,urlcolor=febblue]{hyperref}

\definecolor{febblue}{HTML}{1C4687}
\definecolor{febgold}{HTML}{FDB515}
\definecolor{codebg}{HTML}{F3F4F6}

\lstset{
  basicstyle=\ttfamily\small,
  backgroundcolor=\color{codebg},
  frame=single, framerule=0pt, framesep=4pt,
  xleftmargin=4pt, xrightmargin=4pt,
  breaklines=true, columns=fullflexible, keepspaces=true,
  showstringspaces=false,
}
\newcommand{\file}[1]{\texttt{#1}}
\newcommand{\cmd}[1]{\texttt{#1}}
\newcommand{\rulebox}[1]{\par\noindent\colorbox{codebg}{\parbox{\dimexpr\linewidth-2\fboxsep}{#1}}\par\medskip}

\pagestyle{fancy}
\fancyhf{}
\fancyhead[L]{\color{febblue}\textbf{FEBAUTO Racing}}
\fancyhead[R]{\color{febblue}\nouppercase{\leftmark}}
\fancyfoot[C]{\thepage}
\renewcommand{\headrulewidth}{0.8pt}
\renewcommand{\headrule}{\color{febgold}\hrule width\headwidth height\headrulewidth}
\setcounter{tocdepth}{1}

% \febtitle{subtitle}{part lines}
\newcommand{\febtitle}[2]{%
\begin{titlepage}
  \centering
  \vspace*{2cm}
  \includegraphics[width=5cm]{../../assets/feb_logo_512.png}\\[1.5cm]
  {\color{febblue}\Huge\bfseries FEBAUTO Racing}\\[6pt]
  {\Large #1}\\[2cm]
  \begin{tabular}{rl}#2\end{tabular}\\[3cm]
  Formula Electric at Berkeley\\
  September 2026, platform release v0.1.0\\[6pt]
  \url{https://github.com/FEBAutonomous/feb-racing}\\
  \url{https://febautonomous.github.io/feb-racing/}
\end{titlepage}
\tableofcontents
\newpage}

\setlist{nosep,leftmargin=*}

\begin{document}
\febtitle{Release 2: starting fresh}{%
  \textbf{For} & anyone who has not written driver code yet, with or without the old devkit\\[3pt]
  \textbf{Time} & about 20 minutes, most of it a download\\[3pt]
  \textbf{Result} & the simulator driving, a ROS 2 terminal, rviz in a browser tab\\}

\section{Before you start}

Install these three, once:

\begin{enumerate}
  \item \textbf{Docker.} Mac: Docker Desktop or OrbStack. Windows: Docker Desktop, and in its settings
  turn on \emph{WSL integration} for your Ubuntu. Linux: \cmd{sudo apt install docker.io} and add
  yourself to the \cmd{docker} group, then log out and in.
  \item \textbf{VS Code} with the \emph{Dev Containers} extension. Cursor works the same way.
  \item \textbf{Foxglove} from \url{https://foxglove.dev/download}. Optional, the browser desktop covers most of it.
\end{enumerate}

Windows: everything below happens in a \textbf{WSL 2 Ubuntu terminal}, not PowerShell. If you do not
have one: \cmd{wsl --install -d Ubuntu} in PowerShell, reboot, open ``Ubuntu'' from the Start menu.

\section{Had the old devkit? Delete it}

Release 2 rewrote the repository, so the old clone cannot be updated. It holds nothing you need
if you never wrote code in its \file{stack/} folder. Remove it:

\begin{lstlisting}
rm -rf ~/feb-racing
\end{lstlisting}

If it lives somewhere else, delete it from there. Skip this section if you never had it.

\section{Setup}

\begin{lstlisting}
mkdir -p ~/FEB && cd ~/FEB
git clone git@github.com:FEBAutonomous/feb-racing.git
cd feb-racing
./feb setup --name "Your Name"
\end{lstlisting}

What happens: the launcher pulls the devkit image (a few GB, once), downloads the simulator app for
your OS, and creates \emph{your box}, a container with ROS 2 Humble, the house nodes and a desktop,
with your whole \file{\textasciitilde/FEB} folder mounted inside it. The box is persistent: it survives
reboots, keeps what you \cmd{pip install}, and builds incrementally.

Your name is how your results appear on the scoreboard. Ask Pranav to add it to
\file{submissions.yaml} with your GitHub login.

No \cmd{git@github.com} access yet? Use \cmd{https://github.com/FEBAutonomous/feb-racing.git}
in the clone line; everything else is the same.

\section{First drive}

\begin{lstlisting}
./feb up
\end{lstlisting}

The simulator window opens. Arrow keys drive. Close the window, or Ctrl-C in the terminal, to stop.
Mac, first time only: if macOS refuses to open the simulator, right-click the app in
\file{\textasciitilde/Library/Application Support/feb-sim/app} and choose Open.

To watch the starter driver race instead of driving yourself:

\begin{lstlisting}
./feb up --driver feb_driver
\end{lstlisting}

\section{Look inside: a terminal, rviz in the browser, Foxglove}

\begin{tabularx}{\linewidth}{lX}
\toprule
Command & Gives you \\
\midrule
\cmd{./feb shell} & a terminal inside the box; \cmd{ros2 topic list}, \cmd{ros2 topic echo /scan} \\
\cmd{./feb shell -{}-gui} & the same, plus the box's desktop in a browser tab at \url{http://localhost:6080} \\
\cmd{./feb code} & VS Code opened inside the box, with the ROS extensions \\
\bottomrule
\end{tabularx}

\medskip
\textbf{rviz in the browser.} With \cmd{./feb up} running in one terminal, open a second one:

\begin{lstlisting}
./feb shell --gui
rviz2 &
\end{lstlisting}

A browser tab opens with a grey desktop; rviz appears on it. In rviz: \emph{Fixed Frame} to
\cmd{odom}, then \emph{Add}, \emph{By topic}, \cmd{/scan} (LaserScan) and \cmd{/odom}
(Odometry). Set the LaserScan style to Points, size 4. If a topic shows nothing, set its
\emph{Reliability Policy} to Best Effort. \cmd{rqt} and \cmd{turtlesim} land on the same desktop.

\textbf{Foxglove app.} Open connection, Foxglove WebSocket, \cmd{ws://localhost:8765}. Add a 3D
panel, turn on \cmd{/scan}; the bridge is already running whenever \cmd{feb up} is.

\section{Day to day}

\begin{tabularx}{\linewidth}{lX}
\toprule
Command & Does \\
\midrule
\cmd{./feb update} & pulls the repo, a newer image if there is one, a newer simulator. Run it when told there is a release. \\
\cmd{./feb stop} & stops a run that is stuck \\
\cmd{./feb reset} & throws the box away and makes a new one; your files are untouched \\
\cmd{./feb practice -{}-driver my\_driver} & scores a driver on the qualifying tracks \\
\cmd{./feb submit} & posts clean runs to the scoreboard and your code to the team repo \\
\bottomrule
\end{tabularx}

\medskip
You never re-clone again. \cmd{./feb update} is the whole upgrade path from here on.

Writing a driver, scoring, the dead-man and the car are in \file{member\_guide.pdf}, next to this file.

\section{When something is wrong}

\begin{tabularx}{\linewidth}{lX}
\toprule
It says & Do \\
\midrule
\texttt{Cannot connect to the Docker daemon} & Docker Desktop is not running; open it and wait for the whale icon to settle \\
\texttt{this clone sits directly in /home/...} & the clone must be inside a folder of its own: \cmd{mkdir \textasciitilde/FEB \&\& mv \textasciitilde/feb-racing \textasciitilde/FEB/} \\
\texttt{unauthorized} while pulling the image & you are logged in to \cmd{ghcr.io} with a stale token: \cmd{docker logout ghcr.io}, run setup again \\
the simulator window never appears & \cmd{./feb logs}; on Windows make sure \file{\textasciitilde/FEB} is inside WSL, not on \file{/mnt/c} \\
the browser desktop will not connect & reload the tab; still nothing, \cmd{./feb reset} \\
\bottomrule
\end{tabularx}

\end{document}
